
\uline{d'%
autre part} (IV 2.11) \uline{on a un homomorphisme}
\[
 i_*:K_\bullet(X_s)\longrightarrow K_\bullet(X).
\]
\uline{Ceci posé}, \uline{le composé}
\[
 i^*i_*:K_\bullet(X_s)\longrightarrow K_\bullet(X)\longrightarrow K_\bullet(X_s)
\]
\uline{est nul}.

En effet, \(i:X_s\to X\) est une immersion régulière de codimension \(1\), dont le faisceau conormal \(\mathbb N\) est \emph{libre} de rang \(1\), car \(X_s\) peut se décrire par une équation globale \(f^*(g)=0\) (où \(g\) est une uniformisante de \(A\)). On a donc \(\lambda_{-1}(\mathbb N)=0\), et on conclut grâce à 7.1.

Nous supposerons dans toute la suite \(X\) \emph{noethérien} et \emph{plat sur} \(S\). Appliquons maintenant la suite exacte d'homotopie (IX 1.1) à la situation : \(X\), muni de l'ouvert \(X_t\) et du sous-schéma fermé \(X_s\) complémentaire :
\[
K_\bullet(X_s)\xrightarrow{i_*}K_\bullet(X)\xrightarrow{j^*}K_\bullet(X_t)\longrightarrow0,
\]
où \(j:X_t\to X\) est l'inclusion, qui est une immersion ouverte. Utilisant 7.2, on trouve :

\textbf{Proposition 7.3.} \uline{Avec les notations de} 7.2, \(X\) \uline{étant supposé noethérien}, \uline{il existe un unique homomorphisme de groupes}
\begin{equation}\tag{7.3.1}
 \sigma_\bullet:K_\bullet(X_t)\longrightarrow K_\bullet(X_s)
\end{equation}
(\uline{dit} «~homomorphisme de spécialisation~») \uline{rendant commutatif le diagramme}
\begin{equation}\tag{7.3.2}
\begin{tikzcd}[column sep=large,row sep=large]
K_\bullet(X) \arrow[r,"j^*"] \arrow[d,"i^*"'] & K_\bullet(X_t) \arrow[dl,"\sigma_\bullet"]\\
K_\bullet(X_s) & {.}
\end{tikzcd}
\end{equation}



% Local macro additions for Exposé X continuation
\providecommand{\Kalg}{K_{\mathrm{alg}}}
\providecommand{\Gralg}{\operatorname{Gr}_{\mathrm{alg}}}
\providecommand{\Grnum}{\operatorname{Gr}_{\mathrm{num}}}
\providecommand{\Knum}{K_{\mathrm{num}}}
\providecommand{\Cl}{\operatorname{Cl}}
\providecommand{\Ch}{\operatorname{Ch}}
\providecommand{\Nm}{\operatorname{N}}
\providecommand{\colim}{\varinjlim}
\providecommand{\Hens}{\operatorname{Hens}}
\providecommand{\Spec}{\operatorname{Spec}}
\providecommand{\coker}{\operatorname{coker}}


